La comunicación entre frontend y backend se realiza mediante \textbf{API REST}

\paragraph{Frontend}
Por el lado del frontend, gracias al modulo HttpClient de Angular, se realizan peticiones HTTP al servidor.
Este módulo se utiliza mediante la inyección de dependencias, lo que facilita su uso en los componentes de la aplicación.
Con el uso de este módulo, se han implementado tres servicios principales que se encargan de gestionar las operaciones relacionadas con el backend:
\begin{description}
    \item [AuthService:] se encarga de enviar las peticiones relacionadas con la autenticación, registro y cierre de sesión de los usuarios.
    \item [SearchService:] encargado de las peticiones GET.
    \item [UpdateService:] encargado de las peticiones relacionadas con la actualización o eliminación de los recursos del servidor.
\end{description} 

Una vez que el servidor responde, se transforman los datos en objetos comprensibles por Angular mediante funciones de mapeo.
Dicha transformación permite modifcar los datos para adaptarlos a las necesidades de la aplicación, como por ejemplo,
mostrar el genero del videojuego de forma más legible en la aplicación o aplicar funciones de filtrado sobre los objetos 
transformados.

En la figura \ref{FIG:MAP}, se observa como se realiza el mapeo de los datos a objetos de Angular. 

\begin{figure}[Mapeo de respuesta]{FIG:MAP}{En esta figura se muestra el mapeo de los datos al objeto Videojuegos}
    \image{10cm}{}{videojuegos}
\end{figure}

\paragraph{Backend}
Por el lado del backend, las rutas de la API están protegidas mediante la librería Djoser que valida automáticamente la presencia y validez del token JWT.
En la figura \ref{COD:CODIGOPY} se muestra cómo se protegen las rutas de la aplicación.

\PythonCode[COD:CODIGOPY]{Protección urls}{En esta figura se presenta la protección de las rutas usando Djoser}{python/urls.py}{17}{24}{1}

Si el token, no está presente o no es válido, se devuelve un error 401
En caso contrario, cada viewset ejecuta la lógica correspondiente, implicando operaciones de búsquedas, actualizaciones o inserciones en 
la base de datos mediante el ORM de Django. Por último, la respuesta
se envía al frontend en formato JSON, junto con el código de respuesta HTTP correspondiente. En la figura \ref{FIG:CODIGOPY2} se muestra 
un ejemplo de respuesta del servidor al frontend.

\begin{figure}[Respuesta del servidor]{FIG:CODIGOPY2}{En esta figura se muestra la lista de videojugos devuelta por el servidor}
    \image{10cm}{}{respuesta}
\end{figure}
